
\subsection{Zero order term: Fokker-Planck equation}

The starting point of the numerical calculations for the Fokker-Planck
equation is the following expression 
\begin{equation}
\begin{array}{l}
\frac{\partial f_{0}^{\left( 0\right) }}{\partial t}+\frac{1}{p^{2}}\frac{
\partial }{\partial p}\left( p^{2}\mathbf{S}_{p}^{\left( 0\right) }\right) -%
\frac{1}{\lambda }\frac{1}{p}\frac{\partial }{\partial \xi _{0}}\left( \sqrt{
1-\xi _{0}^{2}}\lambda \mathbf{S}_{\xi }^{\left( 0\right) }\right) + \\ 
\frac{B_{\min }\left( \rho \right) }{\lambda \left( \rho ,\xi _{0}\right) 
\widetilde{q}\left( \rho \right) }\rho _{\psi }^{\prime }\frac{\partial }{
\partial \rho }\left( R_{\min }\left( \rho \right) \widetilde{q}\left( \rho
\right) \frac{B_{P\min }\left( \rho \right) }{B_{\min }\left( \rho \right) }%
\lambda \left( \rho ,\xi _{0}\right) \mathbf{S}_{\rho }^{\left( 0\right)
}\right) =0
\end{array}
\end{equation}

In order to avoid numerical singularities at $p=0$, it is convenient to
multiply all the terms by the partial Jacobian 
\begin{equation}
J_{\psi }J_{p}=\frac{\widetilde{q}\left( \rho \right) }{B_{\min }\left( \rho
\right) }p^{2}
\end{equation}
which leads to the equivalent form 
\begin{equation}
\begin{array}{l}
\frac{\partial \left[ \frac{\widetilde{q}\left( \rho \right) }{B_{\min
}\left( \rho \right) }p^{2}f_{0}^{\left( 0\right) }\right] }{\partial t}+%
\frac{\widetilde{q}\left( \rho \right) }{B_{\min }\left( \rho \right) }\frac{%
\partial }{\partial p}\left( p^{2}\mathbf{S}_{p}^{\left( 0\right) }\right) -%
\frac{\widetilde{q}\left( \rho \right) }{B_{\min }\left( \rho \right) }\frac{
p}{\lambda \left( \rho ,\xi _{0}\right) }\frac{\partial }{\partial \xi _{0}}
\left( \sqrt{1-\xi _{0}^{2}}\lambda \mathbf{S}_{\xi }^{\left( 0\right)
}\right) + \\ 
p^{2}\frac{\rho _{\psi }^{\prime }}{\lambda \left( \rho ,\xi _{0}\right) }%
\frac{\partial }{\partial \rho }\left( R_{\min }\left( \rho \right) 
\widetilde{q}\left( \rho \right) \frac{B_{P\min }\left( \rho \right) }{%
B_{\min }\left( \rho \right) }\lambda \left( \rho ,\xi _{0}\right) \mathbf{S}%
_{\rho }^{\left( 0\right) }\right) =0
\end{array}
\end{equation}

Using the grid definition defined in previous Section **** and applying the
backward time discretization, which corresponds to the fully implicit time
differencing, since the rate of convergence is independent of the time step
even when $\Delta t\gg 1,$ the discrete form of the Fokker-planck equation
may be expressed as

\begin{eqnarray}
&&\frac{\widetilde{q}\left( \rho _{l+1/2}\right) }{B_{\min }\left( \rho
_{l+1/2}\right) }p_{i+1/2}^{2}\frac{f_{0,l+1/2,i+1/2,j+1/2}^{\left( 0\right)
(k+1)}-f_{0,l+1/2,i+1/2,j+1/2}^{\left( 0\right) (k)}}{\Delta t} \\
&&+\frac{\widetilde{q}\left( \rho _{l+1/2}\right) }{B_{\min }\left( \rho
_{l+1/2}\right) }\left. \frac{\partial \left( p^{2}\mathbf{S}_{p}^{\left(
0\right) }\right) }{\partial p}\right| _{l+1/2,i+1/2,j+1/2}^{(k+1)} 
\nonumber \\
&&-\frac{\widetilde{q}\left( \rho _{l+1/2}\right) }{B_{\min }\left( \rho
_{l+1/2}\right) }\frac{p_{i+1/2}}{\lambda \left( \rho _{l+1/2},\xi
_{0,j+1/2}\right) }\left. \frac{\partial }{\partial \xi _{0}}\left( \sqrt{
1-\xi _{0}^{2}}\lambda \left( \rho ,\xi _{0}\right) \mathbf{S}_{\xi
}^{\left( 0\right) }\right) \right| _{l+1/2,i+1/2,j+1/2}^{(k+1)}  \nonumber
\\
&&+p_{i+1/2}^{2}\frac{\rho _{\psi ,_{l+1/2}}^{\prime }}{\lambda \left( \rho
_{l+1/2},\xi _{0,j+1/2}\right) }\left. \frac{\partial }{\partial \rho }%
\left( R_{\min }\left( \rho \right) \widetilde{q}\left( \rho \right) \frac{%
B_{P\min }\left( \rho \right) }{B_{\min }\left( \rho \right) }\lambda \left(
\rho ,\xi _{0}\right) \mathbf{S}_{\rho }^{\left( 0\right) }\right) \right|
_{l+1/2,i+1/2,j+1/2}^{(k+1)}  \nonumber \\
&=&0  \nonumber
\end{eqnarray}
or in a compact form
\begin{eqnarray}
&&\frac{\widetilde{q}\left( \rho _{l+1/2}\right) }{B_{\min }\left( \rho
_{l+1/2}\right) }p_{i+1/2}^{2}\frac{f_{0,l+1/2,i+1/2,j+1/2}^{\left( 0\right)
(k+1)}-f_{0,l+1/2,i+1/2,j+1/2}^{\left( 0\right) (k)}}{\Delta t} \\
&&+\left. \frac{\widetilde{q}\left( \rho \right) }{B_{\min }\left( \rho
\right) }p^{2}\frac{\partial f^{\left( 0\right) }}{\partial p}\right|
_{l+1/2,i+1/2,j+1/2}^{(k+1)}  \nonumber \\
&&+\left. \frac{\widetilde{q}\left( \rho \right) }{B_{\min }\left( \rho
\right) }p^{2}\frac{\partial f^{\left( 0\right) }}{\partial \rho }\right|
_{l+1/2,i+1/2,j+1/2}^{(k+1)}  \nonumber \\
&=&0  \nonumber
\end{eqnarray}
where 
\begin{eqnarray}
\left. \frac{\widetilde{q}\left( \rho \right) }{B_{\min }\left( \rho \right) 
}p^{2}\frac{\partial f^{\left( 0\right) }}{\partial p}\right|
_{l+1/2,i+1/2,j+1/2}^{(k+1)} &=&\frac{\widetilde{q}\left( \rho
_{l+1/2}\right) }{B_{\min }\left( \rho _{l+1/2}\right) }\left. \frac{%
\partial \left( p^{2}\mathbf{S}_{p}^{\left( 0\right) }\right) }{\partial p}%
\right| _{l+1/2,i+1/2,j+1/2}^{(k+1)}  \nonumber \\
&&-\frac{\widetilde{q}\left( \rho _{l+1/2}\right) }{B_{\min }\left( \rho
_{l+1/2}\right) }\frac{p_{i+1/2}}{\lambda \left( \rho _{l+1/2},\xi
_{0,j+1/2}\right) }\left. \frac{\partial }{\partial \xi _{0}}\left( \sqrt{
1-\xi _{0}^{2}}\lambda \left( \rho ,\xi _{0}\right) \mathbf{S}_{\xi
}^{\left( 0\right) }\right) \right| _{l+1/2,i+1/2,j+1/2}^{(k+1)}
\end{eqnarray}
and
\begin{equation}
\left. \frac{\widetilde{q}\left( \rho \right) }{B_{\min }\left( \rho \right) 
}p^{2}\frac{\partial f^{\left( 0\right) }}{\partial \rho }\right|
_{l+1/2,i+1/2,j+1/2}^{(k+1)}=p_{i+1/2}^{2}\frac{\rho _{\psi
,_{l+1/2}}^{\prime }}{\lambda \left( \rho _{l+1/2},\xi _{0,j+1/2}\right) }%
\left. \frac{\partial }{\partial \rho }\left( R_{\min }\left( \rho \right) 
\widetilde{q}\left( \rho \right) \frac{B_{P\min }\left( \rho \right) }{%
B_{\min }\left( \rho \right) }\lambda \left( \rho ,\xi _{0}\right) \mathbf{S}%
_{\rho }^{\left( 0\right) }\right) \right| _{l+1/2,i+1/2,j+1/2}^{(k+1)}
\end{equation}

\subsection{First order term: Drift kinetic equation}

Since the first order drift kinetic equation may be also expressed in a
conservative form as for the Fokker-Planck one, a similar formalisme may be
kept. In that case,
\begin{equation}
\frac{1}{p^{2}}\frac{\partial }{\partial p}\left( p^{2}\mathbf{S}%
_{p}^{\left( 0\right) }\right) -\frac{1}{\lambda }\frac{1}{p}\frac{\partial 
}{\partial \xi _{0}}\left( \sqrt{1-\xi _{0}^{2}}\lambda \mathbf{S}_{\xi
}^{\left( 0\right) }\right) =\frac{1}{p^{2}}\frac{\partial }{\partial p}
\left( p^{2}\widetilde{\mathbf{S}}_{p}^{\left( 0\right) }\right) -\frac{1}{%
\lambda }\frac{1}{p}\frac{\partial }{\partial \xi _{0}}\left( \sqrt{1-\xi
_{0}^{2}}\lambda \widetilde{\mathbf{S}}_{\xi }^{\left( 0\right) }\right) 
\end{equation}
where in the left-hand side, the fluxes concern the function $g^{\left(
0\right) }\left( \psi ,p,\xi _{0}\right) $ which is to be determined,
knowing $\widetilde{\mathbf{S}}_{p}^{\left( 0\right) }$ and  $\widetilde{%
\mathbf{S}}_{\xi }^{\left( 0\right) }$ from the spatial derivative of $%
f_{0}^{\left( 0\right) }\left( \psi ,p,\xi _{0}\right) $given by the zero
order Fokker-Planck equation. Multiplying both sides by the partial Jacobian 
$J_{\psi }J_{p}$, one obtains
\begin{equation}
\frac{\partial }{\partial p}\left( p^{2}\mathbf{S}_{p}^{\left( 0\right)
}\right) -\frac{p}{\lambda }\frac{\partial }{\partial \xi _{0}}\left( \sqrt{
1-\xi _{0}^{2}}\lambda \mathbf{S}_{\xi }^{\left( 0\right) }\right) =\frac{%
\partial }{\partial p}\left( p^{2}\widetilde{\mathbf{S}}_{p}^{\left(
0\right) }\right) -\frac{p}{\lambda }\frac{\partial }{\partial \xi _{0}}
\left( \sqrt{1-\xi _{0}^{2}}\lambda \widetilde{\mathbf{S}}_{\xi }^{\left(
0\right) }\right) 
\end{equation}
and the discrete form of the first order drift kinetic equation is then
simply

\begin{eqnarray}
&&\left. \frac{\partial \left( p^{2}\mathbf{S}_{p}^{\left( 0\right) }\right) 
}{\partial p}\right| _{l+1/2,i+1/2,j+1/2}^{(k)}-\frac{p_{i+1/2}}{\lambda
\left( \rho _{l+1/2},\xi _{0,j+1/2}\right) }\left. \frac{\partial }{\partial 
\xi _{0}}\left( \sqrt{1-\xi _{0}^{2}}\lambda \left( \rho ,\xi _{0}\right) 
\mathbf{S}_{\xi }^{\left( 0\right) }\right) \right|
_{l+1/2,i+1/2,j+1/2}^{(k)}  \nonumber \\
&=&\left. \frac{\partial \left( p^{2}\widetilde{\mathbf{S}}_{p}^{\left(
0\right) }\right) }{\partial p}\right| _{l+1/2,i+1/2,j+1/2}^{(k)}-\frac{
p_{i+1/2}}{\lambda \left( \rho _{l+1/2},\xi _{0,j+1/2}\right) }\left. \frac{%
\partial }{\partial \xi _{0}}\left( \sqrt{1-\xi _{0}^{2}}\lambda \left( \rho
,\xi _{0}\right) \widetilde{\mathbf{S}}_{\xi }^{\left( 0\right) }\right)
\right| _{l+1/2,i+1/2,j+1/2}^{(k)}
\end{eqnarray}
